% Options for packages loaded elsewhere
\PassOptionsToPackage{unicode}{hyperref}
\PassOptionsToPackage{hyphens}{url}
%
\documentclass[
]{article}
\usepackage{amsmath,amssymb}
\usepackage{iftex}
\ifPDFTeX
  \usepackage[T1]{fontenc}
  \usepackage[utf8]{inputenc}
  \usepackage{textcomp} % provide euro and other symbols
\else % if luatex or xetex
  \usepackage{unicode-math} % this also loads fontspec
  \defaultfontfeatures{Scale=MatchLowercase}
  \defaultfontfeatures[\rmfamily]{Ligatures=TeX,Scale=1}
\fi
\usepackage{lmodern}
\ifPDFTeX\else
  % xetex/luatex font selection
\fi
% Use upquote if available, for straight quotes in verbatim environments
\IfFileExists{upquote.sty}{\usepackage{upquote}}{}
\IfFileExists{microtype.sty}{% use microtype if available
  \usepackage[]{microtype}
  \UseMicrotypeSet[protrusion]{basicmath} % disable protrusion for tt fonts
}{}
\makeatletter
\@ifundefined{KOMAClassName}{% if non-KOMA class
  \IfFileExists{parskip.sty}{%
    \usepackage{parskip}
  }{% else
    \setlength{\parindent}{0pt}
    \setlength{\parskip}{6pt plus 2pt minus 1pt}}
}{% if KOMA class
  \KOMAoptions{parskip=half}}
\makeatother
\usepackage{xcolor}
\usepackage[margin=1in]{geometry}
\usepackage{color}
\usepackage{fancyvrb}
\newcommand{\VerbBar}{|}
\newcommand{\VERB}{\Verb[commandchars=\\\{\}]}
\DefineVerbatimEnvironment{Highlighting}{Verbatim}{commandchars=\\\{\}}
% Add ',fontsize=\small' for more characters per line
\usepackage{framed}
\definecolor{shadecolor}{RGB}{248,248,248}
\newenvironment{Shaded}{\begin{snugshade}}{\end{snugshade}}
\newcommand{\AlertTok}[1]{\textcolor[rgb]{0.94,0.16,0.16}{#1}}
\newcommand{\AnnotationTok}[1]{\textcolor[rgb]{0.56,0.35,0.01}{\textbf{\textit{#1}}}}
\newcommand{\AttributeTok}[1]{\textcolor[rgb]{0.13,0.29,0.53}{#1}}
\newcommand{\BaseNTok}[1]{\textcolor[rgb]{0.00,0.00,0.81}{#1}}
\newcommand{\BuiltInTok}[1]{#1}
\newcommand{\CharTok}[1]{\textcolor[rgb]{0.31,0.60,0.02}{#1}}
\newcommand{\CommentTok}[1]{\textcolor[rgb]{0.56,0.35,0.01}{\textit{#1}}}
\newcommand{\CommentVarTok}[1]{\textcolor[rgb]{0.56,0.35,0.01}{\textbf{\textit{#1}}}}
\newcommand{\ConstantTok}[1]{\textcolor[rgb]{0.56,0.35,0.01}{#1}}
\newcommand{\ControlFlowTok}[1]{\textcolor[rgb]{0.13,0.29,0.53}{\textbf{#1}}}
\newcommand{\DataTypeTok}[1]{\textcolor[rgb]{0.13,0.29,0.53}{#1}}
\newcommand{\DecValTok}[1]{\textcolor[rgb]{0.00,0.00,0.81}{#1}}
\newcommand{\DocumentationTok}[1]{\textcolor[rgb]{0.56,0.35,0.01}{\textbf{\textit{#1}}}}
\newcommand{\ErrorTok}[1]{\textcolor[rgb]{0.64,0.00,0.00}{\textbf{#1}}}
\newcommand{\ExtensionTok}[1]{#1}
\newcommand{\FloatTok}[1]{\textcolor[rgb]{0.00,0.00,0.81}{#1}}
\newcommand{\FunctionTok}[1]{\textcolor[rgb]{0.13,0.29,0.53}{\textbf{#1}}}
\newcommand{\ImportTok}[1]{#1}
\newcommand{\InformationTok}[1]{\textcolor[rgb]{0.56,0.35,0.01}{\textbf{\textit{#1}}}}
\newcommand{\KeywordTok}[1]{\textcolor[rgb]{0.13,0.29,0.53}{\textbf{#1}}}
\newcommand{\NormalTok}[1]{#1}
\newcommand{\OperatorTok}[1]{\textcolor[rgb]{0.81,0.36,0.00}{\textbf{#1}}}
\newcommand{\OtherTok}[1]{\textcolor[rgb]{0.56,0.35,0.01}{#1}}
\newcommand{\PreprocessorTok}[1]{\textcolor[rgb]{0.56,0.35,0.01}{\textit{#1}}}
\newcommand{\RegionMarkerTok}[1]{#1}
\newcommand{\SpecialCharTok}[1]{\textcolor[rgb]{0.81,0.36,0.00}{\textbf{#1}}}
\newcommand{\SpecialStringTok}[1]{\textcolor[rgb]{0.31,0.60,0.02}{#1}}
\newcommand{\StringTok}[1]{\textcolor[rgb]{0.31,0.60,0.02}{#1}}
\newcommand{\VariableTok}[1]{\textcolor[rgb]{0.00,0.00,0.00}{#1}}
\newcommand{\VerbatimStringTok}[1]{\textcolor[rgb]{0.31,0.60,0.02}{#1}}
\newcommand{\WarningTok}[1]{\textcolor[rgb]{0.56,0.35,0.01}{\textbf{\textit{#1}}}}
\usepackage{graphicx}
\makeatletter
\def\maxwidth{\ifdim\Gin@nat@width>\linewidth\linewidth\else\Gin@nat@width\fi}
\def\maxheight{\ifdim\Gin@nat@height>\textheight\textheight\else\Gin@nat@height\fi}
\makeatother
% Scale images if necessary, so that they will not overflow the page
% margins by default, and it is still possible to overwrite the defaults
% using explicit options in \includegraphics[width, height, ...]{}
\setkeys{Gin}{width=\maxwidth,height=\maxheight,keepaspectratio}
% Set default figure placement to htbp
\makeatletter
\def\fps@figure{htbp}
\makeatother
\setlength{\emergencystretch}{3em} % prevent overfull lines
\providecommand{\tightlist}{%
  \setlength{\itemsep}{0pt}\setlength{\parskip}{0pt}}
\setcounter{secnumdepth}{-\maxdimen} % remove section numbering
\ifLuaTeX
  \usepackage{selnolig}  % disable illegal ligatures
\fi
\usepackage{bookmark}
\IfFileExists{xurl.sty}{\usepackage{xurl}}{} % add URL line breaks if available
\urlstyle{same}
\hypersetup{
  pdftitle={Tutorial 2},
  hidelinks,
  pdfcreator={LaTeX via pandoc}}

\title{Tutorial 2}
\author{}
\date{\vspace{-2.5em}2024-10-01}

\begin{document}
\maketitle

\subsection{Setting the Working Directory in
R}\label{setting-the-working-directory-in-r}

You can set the working directory in R to specify where files are read
from or saved. Below is the code for setting the working directory on
different operating systems.

\section{Set the working directory}\label{set-the-working-directory}

\section{Replace ``path/to/your/directory'' with your actual directory
path}\label{replace-pathtoyourdirectory-with-your-actual-directory-path}

\section{Windows:}\label{windows}

setwd(``C:/Users/YourName/Documents/ProjectFolder'')

\section{macOS or Linux:}\label{macos-or-linux}

setwd(``/Users/YourName/Documents/ProjectFolder'')

\subsection{Download the Data}\label{download-the-data}

\begin{Shaded}
\begin{Highlighting}[]
\NormalTok{store.df }\OtherTok{\textless{}{-}} \FunctionTok{read.csv}\NormalTok{(}\StringTok{"http://goo.gl/QPDdMl"}\NormalTok{)}
\end{Highlighting}
\end{Shaded}

\subsection{Task 1, lets break down the following
boxplots}\label{task-1-lets-break-down-the-following-boxplots}

This creates a single boxplot for the variable p2sales from the store.df
data frame.

This boxplot represents the weekly sales (p2sales) for all stores
combined, without distinguishing between different stores. For a more
comprehensive analysis, you need to break down the data by meaningful
categories, such as by store.

xlab = ``Weekly sales'': Label for the x-axis, representing the weekly
sales values. ylab = ``P2'': Label for the y-axis, which can be used to
indicate the product or item being analyzed (in this case, ``P2''). main
= ``Weekly sales of P2, All stores'': The title of the boxplot.
horizontal = TRUE: The boxplot is drawn horizontally, with the x-axis
showing the sales values and the y-axis indicating ``P2''.

\begin{Shaded}
\begin{Highlighting}[]
\FunctionTok{boxplot}\NormalTok{(store.df}\SpecialCharTok{$}\NormalTok{p2sales, }\AttributeTok{xlab=}\StringTok{"Weekly sales"}\NormalTok{, }\AttributeTok{ylab=}\StringTok{"P2"}\NormalTok{,}
        \AttributeTok{main=}\StringTok{"Weekly sales of P2, All stores"}\NormalTok{, }\AttributeTok{horizontal=}\ConstantTok{TRUE}\NormalTok{)}
\end{Highlighting}
\end{Shaded}

\includegraphics{Tutorial-2--Part-1_files/figure-latex/unnamed-chunk-2-1.pdf}

This code creates multiple boxplots, one for each store in the dataset
(storeNum).The \textasciitilde{} operator is used here to create
boxplots for p2sales by store. Each boxplot represents the distribution
of weekly sales (p2sales) for a different store.In R, the
\textasciitilde{} can be read as ``by'', meaning the boxplots will be
grouped based on the unique values of storeNum.

Each unique storeNum will generate a separate boxplot.

ylab = ``Store'': Labels the y-axis, indicating that each boxplot
corresponds to a different store. xlab = ``Weekly unit sales'': Labels
the x-axis, which shows the sales data. las = 1: Rotates the y-axis
labels to be horizontal, making the store labels easier to read. main =
``Weekly Sales of P2 by Store'': The title of the plot. horizontal =
TRUE: The boxplots are horizontal, with sales values on the x-axis.

\begin{Shaded}
\begin{Highlighting}[]
\FunctionTok{boxplot}\NormalTok{(store.df}\SpecialCharTok{$}\NormalTok{p2sales }\SpecialCharTok{\textasciitilde{}}\NormalTok{ store.df}\SpecialCharTok{$}\NormalTok{storeNum, }\AttributeTok{horizontal=}\ConstantTok{TRUE}\NormalTok{,}
        \AttributeTok{ylab=}\StringTok{"Store"}\NormalTok{, }\AttributeTok{xlab=}\StringTok{"Weekly unit sales"}\NormalTok{, }\AttributeTok{las=}\DecValTok{1}\NormalTok{,}
        \AttributeTok{main=}\StringTok{"Weekly Sales of P2 by Store"}\NormalTok{)}
\end{Highlighting}
\end{Shaded}

\includegraphics{Tutorial-2--Part-1_files/figure-latex/unnamed-chunk-3-1.pdf}

\subsection{Factor}\label{factor}

Check if store.df\$storeNum is a factor

\begin{Shaded}
\begin{Highlighting}[]
\FunctionTok{is.factor}\NormalTok{(store.df}\SpecialCharTok{$}\NormalTok{storeNum)}
\end{Highlighting}
\end{Shaded}

\begin{verbatim}
## [1] FALSE
\end{verbatim}

\begin{Shaded}
\begin{Highlighting}[]
\FunctionTok{class}\NormalTok{(store.df}\SpecialCharTok{$}\NormalTok{storeNum)}
\end{Highlighting}
\end{Shaded}

\begin{verbatim}
## [1] "integer"
\end{verbatim}

This line of code converts the variable storeNum in the data frame df1
into a factor.

\begin{Shaded}
\begin{Highlighting}[]
\NormalTok{store.df}\SpecialCharTok{$}\NormalTok{storeNum }\OtherTok{\textless{}{-}} \FunctionTok{factor}\NormalTok{(store.df}\SpecialCharTok{$}\NormalTok{storeNum)}
\end{Highlighting}
\end{Shaded}

\begin{Shaded}
\begin{Highlighting}[]
\FunctionTok{class}\NormalTok{(store.df}\SpecialCharTok{$}\NormalTok{storeNum)}
\end{Highlighting}
\end{Shaded}

\begin{verbatim}
## [1] "factor"
\end{verbatim}

\subsection{Task 2: Discuss the use of {[} {]} brackets which contain
conditional row and column expressions in
R.}\label{task-2-discuss-the-use-of-brackets-which-contain-conditional-row-and-column-expressions-in-r.}

This line of R code creates a histogram for the mpce variable in the df5
data frame, but only for the values of mpce that are less than 30,000.

hist(df5\(mpce[df5\)mpce\textless30000{]})

df5\(mpce: This selects the mpce column from the data frame df5
df5\)mpce \textless{} 30000: This condition checks which values in the
mpce column are less than 30,000. It returns a vector of TRUE or FALSE
values. df5\(mpce[df5\)mpce \textless{} 30000{]}: This subsetting
operation uses the condition to extract only the values of mpce that
satisfy the condition (i.e., mpce values less than 30,000). hist(): The
hist() function creates a histogram of the filtered values.

The expression df5\(mpce[df5\)mpce \textless{} 30000{]} only focuses on
filtering rows based on a condition for one specific column, which is
mpce.

\subsection{Task 3: Construct a
histogram}\label{task-3-construct-a-histogram}

This command is creating a histogram of the store.spend column in the
cust.df data frame.

The breaks argument controls how the data is divided into bins
(intervals) in the histogram.
max(cust.df\(store.spend) calculates the maximum value of the store.spend column.
ceiling(max(cust.df\)store.spend)/10) rounds up this maximum value to
the nearest 10. 0:ceiling(\ldots) creates a sequence of numbers from 0
to this rounded maximum. Multiplying by 10 ensures that the bin width is
set to 10

\begin{Shaded}
\begin{Highlighting}[]
\NormalTok{cust.df }\OtherTok{\textless{}{-}} \FunctionTok{read.csv}\NormalTok{(}\StringTok{"http://goo.gl/PmPkaG"}\NormalTok{)}
\end{Highlighting}
\end{Shaded}

\begin{Shaded}
\begin{Highlighting}[]
\FunctionTok{hist}\NormalTok{(cust.df}\SpecialCharTok{$}\NormalTok{store.spend, }
     \AttributeTok{breaks=}\NormalTok{(}\DecValTok{0}\SpecialCharTok{:}\FunctionTok{ceiling}\NormalTok{(}\FunctionTok{max}\NormalTok{(cust.df}\SpecialCharTok{$}\NormalTok{store.spend)}\SpecialCharTok{/}\DecValTok{10}\NormalTok{))}\SpecialCharTok{*}\DecValTok{10}\NormalTok{,}
     \AttributeTok{main=}\StringTok{"Customers as of June 2014"}\NormalTok{, }
     \AttributeTok{xlab=}\StringTok{"Prior 12 months in{-}store sales ($)"}\NormalTok{, }
     \AttributeTok{ylab=}\StringTok{"Count of customers"}\NormalTok{)}
\end{Highlighting}
\end{Shaded}

\includegraphics{Tutorial-2--Part-1_files/figure-latex/unnamed-chunk-9-1.pdf}

\subsection{Task 4:Create a kernel density
diagram}\label{task-4create-a-kernel-density-diagram}

d \textless- density(store.df\(p2prom):
density() generates a kernel density estimate.
store.df\)p2prom is referring to a column named p2prom in the dataframe
store.df. The result is stored in d, containing the estimated density
data. plot(d): This plots the density estimate stored in d.~

\begin{Shaded}
\begin{Highlighting}[]
\NormalTok{d }\OtherTok{\textless{}{-}} \FunctionTok{density}\NormalTok{(store.df}\SpecialCharTok{$}\NormalTok{p2prom)}
\FunctionTok{plot}\NormalTok{(d)}
\end{Highlighting}
\end{Shaded}

\includegraphics{Tutorial-2--Part-1_files/figure-latex/unnamed-chunk-10-1.pdf}

\end{document}
